\section*{Introduction}
\phantomsection\addcontentsline{toc}{section}{Introduction}

An accurate and safe algorithm is only one part of a clinical product. To be used in a
screening programme, DOTS must also be built as medical-device software: its components
must be verifiable, its risks analysed, its decisions traceable and explainable, and it
must integrate with the hospital information system. In Europe, these requirements are
set by the IEC~62304 and ISO~14971 standards and by the EU AI Act, which classifies AI
used in medical devices as high-risk.

This chapter describes how these requirements shaped the design of DOTS from the start,
rather than being added afterwards. We first map each standard and article to the
architectural mechanism that addresses it, then present the verification strategy, the
risk analysis, the latency and scalability of the system, its integration with electronic
health records and its three-tier deployment architecture. We end with the audit trail and
explanation template, the ethical considerations, and an explicit statement of what this
design approach does not solve.

\section{Regulatory Compliance as a Design Principle}
\label{sec:regprinciple}

The DR grading systems reviewed in \cref{ch:sota} treat regulatory compliance as a
post-hoc assessment: a system is built to maximise accuracy, and a regulatory dossier is
assembled afterwards to argue that it is acceptably safe. DOTS inverts this sequence.
Every major design decision of \cref{ch:dots} was motivated, in part, by the requirement
to satisfy IEC~62304~\cite{iec62304}, ISO~14971~\cite{iso14971} and Articles~9--15 of the
EU AI Act~\cite{euaiact2024} by construction, not by retrofit.

The distinction matters in practice. A system designed for accuracy and retrofitted for
compliance typically needs architectural surgery when a gap is discovered late. For
example, it may turn out that no component can produce a per-prediction explanation
satisfying Article~13. A separate post-hoc explainability module (Grad-CAM, SHAP,
LIME~\cite{selvaraju2017gradcam,lundberg2017shap,ribeiro2016lime}) then has to be added,
and its explanations may not faithfully represent the model's actual decision process. In
DOTS, the GAT attention weights $\alpha_{ij}$ (\cref{eq:att}) used to make the prediction
\emph{are} the explanation, not a separate approximation of it.

\section{Mapping Regulatory Standards to Architectural Mechanisms}
\label{sec:regmap}

\Cref{tab:regmap} maps each standard or article to the DOTS mechanism that addresses it.

\begin{table}[H]
\centering
\caption[Regulatory standards mapped to DOTS mechanisms]{Regulatory standards mapped to DOTS mechanisms.}
\label{tab:regmap}
\small
\begin{tabular}{@{}P{3.3cm}P{3.9cm}P{7.2cm}@{}}
\toprule
Standard / article & Requirement & DOTS mechanism \\
\midrule
IEC 62304, Sect.\ 5.1--5.8 & Software life cycle, unit verification & Six typed-interface stages; 47 unit tests; $100\%$ interface coverage \\
ISO 14971 & Risk management (FMEA) & Six hazard classes; typed exception at every stage boundary (\cref{sec:iso14971}) \\
EU AI Act Art.~9 & Risk-management system & FMEA $+$ typed exception handling \\
EU AI Act Art.~10 & Data governance & Five-dataset evaluation; documented provenance (\cref{app:datasets}) \\
EU AI Act Art.~11--12 & Technical documentation, logging & Versioned YAML configurations; structured prediction log \\
EU AI Act Art.~13 & Transparency to users & GAT weights $\alpha_{ij}$ $+$ interpretable TDA features \\
EU AI Act Art.~14 & Human oversight & Entropy abstention $\to$ mandatory specialist review \\
EU AI Act Art.~15 & Accuracy and robustness & Five-benchmark validation $+$ calibration analysis \\
FDA SaMD guidance & Software as a Medical Device & Locked algorithm; 510(k) or De Novo pathway \\
\bottomrule
\end{tabular}
\end{table}

\subsection{IEC 62304: Software Life Cycle and Unit Verification}
\label{sec:iec62304}

IEC~62304 classifies medical-device software into three safety classes (A, B, C)
according to the severity of harm that a software failure could cause. DOTS, a
decision-support system whose output informs referral decisions, is treated as Class~B
(non-serious injury possible) under the conservative reading adopted in this thesis,
pending formal classification by a notified body during the Phase~4 submission
(\cref{sec:roadmap}).

Class~B requires documented unit verification. The six-stage pipeline
(\cref{tab:stages}) was designed for this purpose: each stage has a typed input/output
contract (for example, stage S3 accepts a skeleton point cloud $P_i\subset\mathbb{R}^2$
and returns $x_i^{\mathrm{TDA}}\in\mathbb{R}^7$, with explicit shape and dtype assertions
at the function boundary), so that each stage can be tested in isolation.
\Cref{tab:unittests} gives the number of tests and the interface coverage of each stage.

\begin{table}[H]
\centering
\caption[Unit-test coverage by pipeline stage (IEC 62304 verification)]{Unit-test coverage by pipeline stage (IEC 62304 verification).}
\label{tab:unittests}
\small
\begin{tabular}{@{}c c P{8.4cm} c@{}}
\toprule
Stage & Tests & Coverage focus & Interface coverage \\
\midrule
S1 & 9 & Image dimension and dtype, CLAHE parameter bounds, skeleton connectivity & 100\% \\
S2 & 6 & Embedding dimension, NaN/Inf detection, batch consistency & 100\% \\
S3 & 11 & Persistence-diagram validity, empty-skeleton edge case, filtration bounds & 100\% \\
S4 & 4 & Concatenation dimension, type consistency & 100\% \\
S5 & 8 & Graph connectivity, $k$-NN correctness, edge weights in $[0,1]$ & 100\% \\
S6 & 9 & CORAL monotonicity (property-based), entropy bounds, abstention logic & 100\% \\
\midrule
Total & 47 & & 100\% \\
\bottomrule
\end{tabular}
\end{table}

The stage-S6 monotonicity test deserves a note. Rather than testing specific
input/output pairs, it is a property-based test that generates $10{,}000$ random
embeddings $z'_i$ and asserts $\hat P(Y\geq1)\geq\hat P(Y\geq2)\geq\cdots$ for each of
them. It verifies the algebraic guarantee of \cref{eq:coral-proj} at unit-test level,
independently of any trained weights.

\subsection{ISO 14971: Risk Management and FMEA}
\label{sec:iso14971}

ISO~14971 requires the systematic identification and mitigation of hazards. The Failure
Mode and Effects Analysis (FMEA) of DOTS identifies six hazards, each with a severity class
(Critical, Major, Moderate), a pipeline stage of origin and an architectural mitigation
(\cref{tab:fmea}). Three hazards are Critical, two are Major and one is Moderate.

\begingroup
\small\setstretch{1.0}
\setlength{\tabcolsep}{4pt}
\begin{longtable}{@{}c l P{4.3cm} P{3.4cm} P{5.4cm}@{}}
\caption[Complete FMEA across the six identified hazards]{Complete FMEA across the six identified hazards.}\label{tab:fmea}\\
\toprule
\# & Severity & Failure mode & Potential effect & Mitigation \\
\midrule
\endfirsthead
\multicolumn{5}{@{}l}{\small\itshape \Cref{tab:fmea}, continued}\\
\toprule
\# & Severity & Failure mode & Potential effect & Mitigation \\
\midrule
\endhead
\bottomrule
\endlastfoot
1 & Critical & Corrupted, ungradable or wrong-modality image at stage S1 & Garbage-in, garbage-out prediction with false confidence & Typed input contract rejects non-conforming images at S1; the image is flagged for retake instead of being processed silently \\
2 & Major & Skeleton extraction fails on a low-contrast or poorly illuminated image (S1) & Degraded S3 features; possible silent accuracy loss & CLAHE reduces incidence; a typed exception is raised on extraction failure instead of returning a degenerate vector \\
3 & Major & Backbone or TDA feature drift from the training distribution (S2--S3), e.g.\ a new camera or population & Reduced accuracy, possibly undetected & Distribution-shift monitoring flags out-of-envelope feature vectors; flagged predictions are routed to human review \\
4 & Moderate & Insufficient graph neighbours at a new site with a small initial cohort (S5) & Weaker relational signal during programme ramp-up & Minimum-$k$ fallback; explicit low-confidence flag when fewer than five neighbours are available \\
5 & Critical & Over-confident prediction near a grade boundary (S6) leading to an unsafe automated decision & Missed urgent referral (severe DR or PDR) & CORAL rank guarantee (\cref{eq:coral-proj}) combined with entropy abstention (\cref{sec:abst}); the hazard targeted by the H3a safety criteria (\cref{tab:h3a}) \\
6 & Critical & Prediction not delivered, misrouted or silently dropped during EHR integration (S6 $\to$ Tier~3) & Actionable prediction never reaches the treating clinician & FHIR R4 adapter with delivery confirmation (\cref{sec:fhir}); a Phase~2 deliverable, so the hazard is only partially mitigated at the time of writing \\
\end{longtable}
\endgroup

\paragraph{Residual risk acceptance.}
Hazard~6 is the only Critical hazard whose mitigation is incomplete. Its residual risk is
acceptable for the Level~1 computational validation of this thesis, but not for clinical
deployment without the Phase~2 FHIR adapter; this distinction is the basis of the phased
roadmap (\cref{sec:roadmap}), not an oversight.

\paragraph{Severity definitions.}
\emph{Critical}: potential for irreversible patient harm (missed sight-threatening
diagnosis). \emph{Major}: potential for delayed but recoverable harm, or systematic
accuracy loss requiring intervention. \emph{Moderate}: workflow disruption without direct
patient-safety implications.

Each Critical hazard has a documented architectural mitigation: input quality is handled
by typed contract rejection at stage S1; over-confident boundary predictions by the CORAL
guarantee combined with entropy abstention. EHR integration failure is the one Critical
hazard whose mitigation (FHIR R4 delivery confirmation, \cref{sec:fhir}) is still a
Phase~2 deliverable rather than a present capability.

\section{EU AI Act Annex III: High-Risk Classification}
\label{sec:aiact}

Annex~III of the EU AI Act classifies AI systems used as safety components of medical
devices as high-risk, which triggers the requirements of Articles~9--15 before CE
marking. DOTS addresses each of them by construction.
\begin{description}[style=nextline,leftmargin=1.2em]
\item[Art.~9 (risk management)] The FMEA of \cref{tab:fmea}, combined with typed
exception handling at every stage boundary, forms a documented, maintainable
risk-management system.
\item[Art.~10 (data governance)] Five-dataset evaluation (\cref{app:datasets}) with
documented provenance, licence terms and known limitations (for example the small size
of IDRiD).
\item[Art.~11--12 (technical documentation, logging)] Versioned YAML configurations for
all twelve principal experiments (\cref{app:results}) and the structured prediction log
of \cref{sec:audit}.
\item[Art.~13 (transparency)] GAT attention weights and TDA feature values, surfaced
through the explanation template of \cref{sec:audit}, give transparency intrinsic to the
inference path rather than a post-hoc approximation.
\item[Art.~14 (human oversight)] Entropy-based abstention (\cref{sec:abst}) makes human
oversight mandatory for exactly the cases where the model is least reliable.
\item[Art.~15 (accuracy, robustness)] The five-benchmark validation (\cref{tab:bench}) and
the calibration analysis (\cref{tab:ece}) document accuracy and robustness across
acquisition conditions (four countries, several camera models).
\end{description}

\section{Scalability and Per-Patient Latency}
\label{sec:latency}

\Cref{tab:latency} breaks the per-patient latency down by stage. Stage S3 (TDA) accounts
for $57\%$ of the total and is the principal computational bottleneck; GPU-accelerated
persistent homology (for example with giotto-tda~\cite{tauzin2021giotto}) could in
principle reduce it substantially, but this has not been validated within the thesis.

\begin{table}[!tb]
\centering
\caption[Per-patient latency by stage]{Per-patient latency by stage. The cost of S6 is $O(S^L)$, independent of cohort
size.}
\label{tab:latency}
\small
\begin{tabularx}{\textwidth}{@{}l X l@{}}
\toprule
Stage & Complexity & Latency per patient \\
\midrule
S1 -- Pre-processing & $O(W\times H)$ & 4.2~ms \\
S2 -- CNN/ViT extraction & $O(1)$ per image & 6.8~ms \\
S3 -- TDA (Rips filtration) & $O(n\log n)$, $n$ skeleton points & 18.0~ms (CPU, 16 cores) \\
S4 -- Feature fusion & $O(1)$ & $<0.1$ ms \\
S5 -- Graph construction (FAISS) & $O(k\log N)$ per query & 1.2~ms \\
S6 -- GAT/GraphSAGE ($S=25$, $L=2$) & $O(S^L)=O(625)$ & 1.2~ms \\
\midrule
Total & & 31.4~ms \\
\bottomrule
\end{tabularx}
\tabnote{At 31.4~ms per image, one GPU processes about 115{,}000 images per hour when
images are processed one at a time; batching would increase throughput further.}
\end{table}

\paragraph{Throughput at national scale.}
At $31.4$~ms per image, one GPU processes about $115{,}000$ images per hour. A national
programme screening $500{,}000$ patients a year (about $1{,}370$ per day at steady state)
therefore needs a small fraction of one GPU: compute cost is not a barrier to national
deployment.

\paragraph{Memory footprint.}
The population graph must keep the fused embedding of every reference patient. Its $1031$
dimensions ($1024$ from RETFound and $7$ topological features), stored as 32-bit
floating-point numbers, occupy about $4$~kB, so that a reference cohort of $100{,}000$
patients requires about $0.4$~GB, well
within the memory of a single server. The FAISS index adds a comparable amount. The
cohort can therefore grow over the years of a screening programme without changing the
hardware, and new reference patients can be added to the index without retraining the
model.

\section{FHIR R4 Integration Gap}
\label{sec:fhir}

FHIR R4~\cite{hl7fhir} is the international standard for exchanging electronic health
record (EHR) data. Without a FHIR adapter, DOTS outputs cannot be written automatically
to the patient record and would require manual transcription -- a workflow friction that
would likely deter routine clinical adoption regardless of accuracy. Full integration
requires three components.
\begin{enumerate}
\item \textbf{Output mapping.} The DOTS output (grade, confidence vector, entropy,
abstention flag) is mapped to a FHIR \texttt{Observation} resource using the appropriate
SNOMED~CT concept for each ICDR grade.
\item \textbf{Input ingestion.} A DICOM-to-DOTS adapter extracts the fundus image and its
metadata (pseudonymised patient identifier, device, timestamp) from the ophthalmic imaging
workflow.
\item \textbf{Point-of-care embedding.} CDS Hooks integration lets predictions appear
inside the clinician's existing EHR interface rather than in a separate application.
\end{enumerate}
The FHIR adapter is the targeted Phase~2 deliverable (\cref{sec:roadmap}). At the time of
writing it is the largest gap between DOTS as a validated research system and DOTS as a
deployable clinical product.

\section{Three-Tier Deployment Architecture}
\label{sec:tiers}

\Cref{tab:tiers} describes the three deployment tiers. The separation is itself a
regulatory design choice: because Tier~2 (the validated inference engine) is decoupled
from Tier~3 (site-specific integration), a new Tier~3 implementation at a new site does
not require re-validation of Tier~2, provided the Tier~2 software version is unchanged.
This property is directly relevant to the ``locked'' versus ``continuously learning''
distinction in FDA guidance~\cite{fda2017samd}.

\begin{table}[H]
\centering
\caption[Three-tier deployment architecture for national screening]{Three-tier deployment architecture for national screening.}
\label{tab:tiers}
\small
\begin{tabular}{@{}l P{7.4cm} P{4.4cm}@{}}
\toprule
Tier & Function & Status \\
\midrule
1: Acquisition & Fundus camera, DICOM capture, local quality check (S1 rejects ungradable images) & Mature (existing clinical infrastructure) \\
2: Inference & DOTS six-stage pipeline (\cref{tab:stages}) as a containerised inference service & Complete, validated in this thesis \\
3: Integration & FHIR R4 adapter, CDS Hooks, EHR write-back, audit-log persistence & Phase 2 deliverable (\cref{sec:fhir}) \\
\bottomrule
\end{tabular}
\end{table}

\section{Audit Trail and Article 13 Explanation Template}
\label{sec:audit}

Every prediction is logged with the patient pseudonym, the image DICOM UID, a timestamp,
the full five-grade probability vector, the entropy value, the abstention flag and the
GAT attention weights of the five most-attended neighbours. Logs are retained for ten
years, which meets IEC~62304 traceability needs and is compatible with the
storage-limitation principle of GDPR Article~5(1)(e)~\cite{gdpr2016} for medical
records, subject to national retention rules.

The Article~13 explanation, generated automatically from logged inference-time
quantities, reads:

\begin{quote}\small\itshape
Grade [$k$] predicted with [$p$]\% confidence. The five most similar neighbours had
grades [$g_1,\ldots,g_5$] with attention weights [$\alpha_1,\ldots,\alpha_5$]. $H_1$
persistence [$v$] is consistent with Grade [$k$] according to the validated association
of \cref{tab:anova}. Predictive entropy $H_i=[h]$ nats is [below/above] the abstention
threshold $\tau_U=0.73$.
\end{quote}

The template is built entirely from quantities computed during the forward pass. No
separate explanation model is needed, which removes the risk, well documented in the
explainable-AI literature, that a post-hoc explanation misrepresents the actual decision
process.

\subsection{Future Extension: Evidence-Grounded LLM Explanations}
\label{sec:llm}

The fixed template is deliberately simple. A natural extension, not implemented in this
thesis, is to turn the quantities already logged by DOTS into an evidence-grounded
narrative adapted to its reader, clinician or patient. \Cref{fig:llm} presents the
architecture proposed for this purpose. It is organised in eight modules.

\begin{enumerate}[leftmargin=1.6em,itemsep=2pt]
\item \textbf{Input layer.} The fundus image and the clinical data of the patient (age,
sex, diabetes duration, HbA1c, blood pressure, comorbidities).
\item \textbf{Diagnosis engine.} DOTS produces the grade, its confidence, the attention
maps over the image and the graph, and the image and graph embeddings.
\item \textbf{Retrieval module.} A retrieval-augmented generation (RAG) component queries a
vector database of grading criteria, clinical guidelines, textbooks and research papers,
and returns the $k$ most relevant passages for the clinical question.
\item \textbf{Clinical reasoning orchestrator.} A graph is built from the embeddings, the
prediction, the clinical features and the quality of the retrieved evidence; a graph
network with a scoring layer selects the most suitable reasoning strategy among a
library of prompts: strict grounding, evidence first, uncertainty-aware, guideline
comparison, clinician explanation or patient explanation.
\item \textbf{LLM reasoning module.} A large language model uses the selected prompt and
the retrieved evidence to write the explanation and a management plan.
\item \textbf{Output layer.} Grade, key findings, supporting references, management plan,
confidence and uncertainty, phrased for the clinician or for the patient.
\item \textbf{Experience memory.} Every case is stored with its embeddings, evidence,
selected strategy, outcome and feedback, so that similar past cases can guide future
selections.
\item \textbf{Feedback and continuous learning.} Clinician feedback on each explanation is
used to update the selection policy and the memory.
\end{enumerate}

The orchestrator reuses the idea at the heart of this thesis: just as DOTS uses similar
patients to grade a new one, the orchestrator uses similar past cases to choose how to
explain it. For example, when the entropy of the prediction is high, the
uncertainty-aware strategy would state explicitly that the case was ambiguous and why it
was deferred, rather than presenting a confident narrative.

% Figure 6.1 --- patient-aware graph-driven prompt orchestration (future work)
\afterpage{%
\begin{landscape}
\thispagestyle{plain}
\begin{figure}[p]
\centering
\includegraphics[width=\linewidth,height=0.9\textheight,keepaspectratio]{fig_llm_orchestration}
\caption[Patient-aware graph-driven prompt orchestration (future work)]{Patient-aware
graph-driven prompt orchestration for explainable, retrieval-augmented clinical decision
support in diabetic retinopathy (future work, not implemented). The eight modules are
described in \cref{sec:llm}.}
\label{fig:llm}
\end{figure}
\end{landscape}}


Such an extension raises questions outside the scope of this thesis, chiefly how to
prevent a fluent narrative from overriding or contradicting the CORAL ranking and the
abstention decision, which are the load-bearing safety mechanisms. An unfaithful
narrative would undermine exactly the transparency this chapter was designed to provide.
No LLM experiments were conducted; this subsection is a scoped direction for future work.

Should such a layer be built, it would need its own evaluation criteria, distinct from
those of the grading model: \emph{faithfulness} (every statement in the narrative must be
traceable to a logged quantity), \emph{consistency} (the narrative must never contradict the
grade, the confidence or the abstention decision), and \emph{usefulness}, measured with
clinicians and patients in the reader study of Phase~2. These criteria could be
pre-specified in the same way as H1--H3.

\section{Ethical Considerations}
\label{sec:ethics}

Three ethical considerations go beyond formal regulatory requirements.

\paragraph{Algorithmic fairness across demographic subgroups.}
The five benchmarks (\cref{app:datasets}) come from the United States, India, France and
China, which gives some geographic and probably some ethnic diversity. None of them,
however, provides per-image demographic metadata (age, sex, ethnicity), which prevents a
formal subgroup fairness analysis. This is stated as a limitation (\cref{sec:limits-gc})
rather than elided.

\paragraph{Deskilling risk.}
Reliance on automated screening may, over time, erode human grading skill. By routing
$6.2\%$ of cases -- the most challenging ones -- to human review, entropy abstention keeps
graders actively practising on difficult cases rather than removing them from the
workflow.

\paragraph{Accountability for missed referrals.}
Even at Sev-NR $=1.9\%$, some severe cases will be missed. The audit trail of
\cref{sec:audit} makes every missed case traceable to a logged prediction with its full
explanation, which supports clinical-governance review. It does not settle the broader
question of legal and professional accountability when an automated system contributes
to a missed diagnosis; that question is left to the regulatory and legal review of
Phases~3--4 (\cref{sec:roadmap}). In all cases, DOTS is designed as a decision-support tool: the responsibility for the final referral decision remains with the clinician.

\section{What Regulatory-by-Construction Does Not Solve}
\label{sec:regnotsolve}

Designing for compliance from the outset avoids a specific class of late architectural
surgery; it is worth being equally explicit about what it does not solve. Mapping DOTS to
IEC~62304, ISO~14971 and the EU AI Act (\cref{tab:regmap}) documents
\emph{architectural} compliance: the system has the structural features (typed
interfaces, FMEA, audit logging, abstention) that a notified body would expect. It does
not constitute the conformity assessment, clinical evaluation or notified-body review
required before CE marking; those remain Phase~2--4 activities outside the scope of a
doctoral thesis. Likewise, the FHIR gap of \cref{sec:fhir} is not a hypothetical risk: it
is the concrete reason why DOTS could not be connected to a live clinical workflow today.

The honest claim is therefore narrower than ``DOTS is compliant'': DOTS's architecture was
shaped by the requirements it will eventually need to satisfy. This has a practical
consequence for replication. Because each compliance property is embedded in the pipeline
rather than asserted afterwards, any laboratory that reproduces the six-stage pipeline
with the released code (\cref{app:repro}) also reproduces these properties; the FHIR
adapter is the only element that requires site-specific work before clinical use. The
results of \cref{ch:toporet,ch:dots} constitute Level~1 (computational) evidence
in the roadmap of the General Conclusion; regulatory clearance (Level~4) requires
prospective validation, notified-body assessment and post-market surveillance, none of
which is conducted here.

\section{Towards Clinical Deployment: Translation Roadmap}
\label{sec:roadmap}\label{app:roadmap}

The regulatory-oriented design of this chapter prepares DOTS for certification, but does
not replace the clinical evidence that certification requires. The translation of DOTS
into clinical practice is therefore planned in four phases (\cref{tab:roadmap}), following
the four levels of evidence used in diagnostic AI evaluation: Level~1 computational
evidence, Level~2 prospective agreement with human readers, Level~3 randomised patient
outcomes, and Level~4 regulatory clearance. This thesis provides Level~1 evidence only.

\begin{table}[!tb]
\centering
\caption[Four-phase clinical translation roadmap]{Four-phase clinical translation roadmap.}
\label{tab:roadmap}
\small
\begin{tabular}{@{}c l P{7.6cm} l l@{}}
\toprule
Phase & Period & Key deliverable & Evidence & Status \\
\midrule
1 & 2023--2026 & Nine publications; five-benchmark validation; necessity proof & Level 1 & Complete \\
2 & 2026--2027 & Multi-site reader study ($\geq3$ sites, $N\geq200$, pre-registered); FHIR R4 adapter & Level 2 & Planned \\
3 & 2027--2028 & Randomised trial: AI-assisted versus standard-of-care screening & Level 3 & Planned \\
4 & 2028--2029 & CE marking (EU AI Act, high-risk); FDA 510(k) or De Novo & Level 4 & Planned \\
\bottomrule
\end{tabular}
\end{table}

Phase~2 is the critical next step. It combines a pre-registered, multi-site reader study
($N\geq200$ patients, at least three sites), in which board-certified ophthalmologists
grade images with and without DOTS, with the implementation of the FHIR~R4 adapter of
\cref{sec:fhir}. Its primary success criterion is a Sev-NR below $2\%$ with the upper
bound of its 95\% confidence interval also below $2\%$, which would close the main
statistical limitation of the present results. Ethics approval and the FHIR adapter have no
technical dependency and can proceed in parallel. The dates of Phases~2--4 are indicative:
they depend on clinical and industrial partnerships beyond the doctoral work.

\section*{Conclusion}
\phantomsection\addcontentsline{toc}{section}{Conclusion}

This chapter mapped DOTS to IEC~62304, ISO~14971 and the EU AI Act by construction
(\cref{tab:regmap}), documented a 47-test verification suite with full interface coverage
of the six stages, and identified the FHIR R4 adapter (Tier~3) as the main remaining
barrier between a computationally validated system and a clinically deployed product.
Per-patient inference ($31.4$~ms) does not depend on cohort size, so the inductive design
choice of \cref{ch:found} carries through to deployment scale. The three-tier
architecture decouples the validated inference engine from site-specific integration.
Three ethical considerations -- fairness, deskilling and accountability -- are identified
for explicit future work. The General Conclusion now brings together the contributions of the thesis, their
practical value, their limits and the perspectives they open.
